\chapter{Materials and Methods}
\label{sec:methods}


\section{Materials}

\subsection{\acs{C2S} and \acs{C3S} clinker}

Commercially available alite (MIII polymorph, Vustah, Czech Republic) and belite (Vustah, Czech Republic) were utilized.
The chemical composition (Table~\ref{tab:xrf}) was determined by \ac{mXRF}, and the phase purity was confirmed at \num{99.4}\,±\,\SI{0.4}{\wtpercent} (alite) using \ac{XRD} analysis.

\begin{table}[h]
	\centering
	\caption{Chemical composition of the utilized alite obtained by \ac{mXRF}. \label{tab:xrf}}
	\begin{tabular}{@{}lr@{}}
		\toprule
		Oxides  & Alite in \si{\wtpercent}   \\
		\midrule
		SiO$_2$		& 26.34 \\
		TiO$_2$		& 0.031 \\
		Al$_2$O$_3$	& 0.238 \\
		Fe$_2$O$_3$	& 0.094 \\
		Mn$_2$O$_3$	& 0.006 \\
		MgO			& 1.700 \\
		CaO			& 70.79 \\
		Na$_2$O		& 0.083 \\
		\bottomrule
	\end{tabular}
\end{table}

The clinker pastes were prepared with a \ac{wb}-ratio of \num{0.5} and stored in sealed cylindrical containers with a diameter of approx. \SI{8}{\milli\meter}.
%The phase composition and alite purity of similarly prepared samples stored for 1, 7, and 28 days were assessed by quantitative XRD analysis using the G-factor method \cite{jansen.2011}.
%The measurements were carried out using a \textsc{Bruker D8 DaVinci} diffractometer (MA, USA).
%After the required age, the hydration of the prisms was stopped by immersion in isopropanol and subsequent drying at 60$^\circ$C.

\subsection{Porous glass standard}

As reference  with a defined pore size, porous glass membranes by \textit{Boraglas} with a pore size of \num{5} and \SI{20}{\nano\meter} were used.
These borosilicate glass products are made using the \textsc{VYCOR}-process and are usually utilized as filter material.
The specimens were provided as \num{10}\,$\times$\,\num{10}\,$\times$\,\SI{1.2}{\milli\meter} plates, which were cut into smaller pieces for the measurements.

The manufacturer provided a certificate for the porosity of the materials based on \ac{MIP} measurements (Figure~\ref{fig:vycor-mip}).
The \SI{5}{\nano\meter} glass sieve has a modal pore diameter of \num{4.5} and a median diameter of \SI{4.7}{\nano\meter} at a porosity of \SI{54.1}{\percent}.
For the \SI{20}{\nano\meter} material a modal pore diameter of \num{22.3} and a median diameter of \SI{19.4}{\nano\meter} at a porosity of \SI{51.3}{\percent} is given.
A revalidation of the \ac{MIP} results was omitted, due to the high material costs.

\section{The preparation process for SEM imaging}

To obtain high-quality images using a \ac{SEM}, a lot of preparation methods are possible depending on the specimen type.
A selection of methods used within this work are listed here:

\begin{itemize}
	\item \textbf{Scatter sample}. This means sprinkling powdery material over an adhesive carbon pad, followed by the removal of loose material by using compressed air. This method is fairly dangerous for the \ac{SEM} if not done properly.
	\item \textbf{Fractured surface}. Crack the specimen, lean it with compressed air and observe the freshly prepared surface.
	\item \textbf{Cross section (large scale)}. Most mineral specimens have to be embedded in epoxy resin before they are mechanically polished. (Section~\ref{sec:mech_pol}) The polishing result can then be improved by applying \ac{ArBIB}. (Section~\ref{sec:bib_pol_moving})
	\item \textbf{Cross section (small scale)}. A small rectangular specimen (embedded or resin free) can be polished using an \ac{ArBIB} to polish an area of about \SI{1}{\square\milli\meter}. (Section~\ref{sec:bib_pol_fixed})
	\item \textbf{\Ac{FIB} cross section}. Cutting and polishing a small surface (up to 40\,$\times$\,\SI{10}{\square\micro\metre}) using \ac{FIB}. (Section~\ref{sec:fib})
	\item \textbf{Cryo preparation}. To observe water containing specimens (slurries, pastes or very young specimens), the specimen is frozen within \SI{400}{\milli\second} down to about \SI{113}{\kelvin} at \SI{>\,200}{\mega\pascal} to avoid ice-crystallization. The specimen can then be processed by \ac{FIB} or by sublimation.
\end{itemize}

A selection of these methods is described in more detail in the following sections.

\subsection{Mechanical polishing}
\label{sec:mech_pol}

%The standard polished cross section is often used to obtain \ac{BSE} images.
To do a quantitative image analysis, very smooth specimen surfaces are required.
Those surfaces are usually achieved by mechanically polishing epoxy resin embedded specimens.

Before embedding the specimen, the reaction of young specimens has to be stopped first by immersing them into isopropanol to remove the water.
Nevertheless, building materials have to be dried before the embedding process.
Depending on the sensitivity of the specimen, they will remain at \SIrange{40}{60}{\celsius} in an oven for \SI{12}{\hour}.
The sample, which has been brought to the correct size, is now poured over with epoxy resin of a very low viscosity in a mold.
This is required to deeply penetrate the specimen and fill the majority of pores.
In this work, a combination of the resin R1370 and the hardener R1376, by Agar Scientific was used.
The embedded specimen will then be degassed in two \SI{30}{\min} cycles.

After hardening for \SI{24}{\hour} in an oven at \SI{40}{\celsius}, the specimen can be polished.
This is done in multiple steps using diamond-oil-slurries.
The particle sizes of the diamond pastes will be sequentially reduced (\num{15}, \num{9}, \num{3}, \num{1} down to \SI{0.25}{\micro\meter}).
While this polishing method is good enough for many applications like BSE- or EDX-imaging, high-resolution imaging or \ac{EBSD} require even better surface qualities.

\subsection{Argon Broad Ion Beam polishing (moving specimen)}
\label{sec:bib_pol_moving}

Most \ac{ArBIB} devices provide two modes of operation.
In one of the modes, the mechanically polished specimen is semi-erratically moved under a grazing incidence angle under the ion beam.
This method can be applied to improve or clean mechanically polished specimens or to remove thin coatings (e.g. carbon or gold) applied for previous measurements.
In this study, a triple ion beam milling system (EM TIC 3X, Leica) was used (Figure~\ref{fig:bib_device}).
The \ac{ArBIB} polishing is carried out under low pressure (\SIrange{1}{4}{\kilo\pascal}) conditions in an argon atmosphere.

\subsection{Argon Broad Ion Beam sectioning (fixed specimen)}
\label{sec:bib_pol_fixed}

Another mode requires a smaller pre-cut sample (5\,$\times$\,3\,$\times$\,\SI{10}{\cubic\milli\meter}) in a fixed sample holder.
The three argon ion beams intersect at a sharp edged tungsten mask forming a horizontal milling sector of approximately of \SI{100}{\degree} (Figure~\ref{fig:bib_device}).
The mask and the specimen surface are arranged in a way that only \SIrange{20}{100}{\micro\meter} of the sample are exposed to the ion beam above the mask (Figure~\ref{fig:bib_scheme}).

\begin{figure}[!ht]
    \centering

	\begin{subfigure}{0.4\textwidth}
		\includegraphics[width=\textwidth]{methods/bib_device.pdf}
		\caption{Photo of the specimen camber of the \ac{ArBIB} polishing device.}
		\label{fig:bib_device}
	\end{subfigure}
	\hfill
	\begin{subfigure}{0.58\textwidth}
		\includegraphics[width=\textwidth]{methods/bib_scheme.pdf}
		\caption{Principle of \ac{ArBIB} sectioning of a stationary specimen using a tungsten blade, side view. Image based on \cite{kleiner.2021a}}
		\label{fig:bib_scheme}
	\end{subfigure}
    \caption{Illustration of the \ac{ArBIB} sectioning principle for a stationary specimen.}
    \label{fig:bib}
\end{figure}

% bad paper about ArBIB \cite{sammer.2024}

\FloatBarrier
\section{SEM / Focussed Ion Beam}
\label{sec:sem}

In this thesis, mainly a \textit{thermoFischer Scientific Helios G4 UX} was used for the \ac{SEM} and \ac{FIB} measurements.
Some of the \ac{SEM} imaging was done utilizing a \textit{Nova NanoSEM 230} by \textit{FEI}.

These \ac{SEM}-devices can be equipped with a wide range of detectors, stages and tools allowing a wide range of applications.
The most common detectors are \ac{SE} and \ac{BSE} detectors of various types (e.g. \ac{ETD} or \ac{TLD}).
Most modern SEMs easily allow to image large areas using these detectors in a high resolution by acquiring multiple images and stitching them afterwards.


For chemical analysis, \ac{EDX} detectors can provide point and area mappings (see Section~\ref{sec:edx}).


A less common detector type, the \ac{EBSD}, allows to obtain crystallographic information of a clean specimen surface (see Section~\ref{sec:ebsd}).

Besides these classical detectors, modern SEMs can also be combined with a \ac{FIB}.
To apply this method, the cementitious specimen was embedded in resin to avoid imaging the pore background or loose moving material within larger pores.
Resin filled pores have advantages for later pore segmentation.
The specimen should be a polished section to find a suitable \ac{ROI} using the \ac{BSE} detector.

While the \ac{SEM} is widley used to analyze cement based materials, it has its pitfalls.
In contrast to an light microscope, it can alter the specimen during the observation significantly.
%For example, early publications define \ac{CSH} as gel.
%In one part, this due to its low crystallinity, which is hardly detectable using \ac{XRD} but only \ac{TEM} \cite{melzer.1989, rossler.2006}.
%On the other hand, the \ac{CSH}-crystals appeared gel like due to preparation artifacts and severe electron beam damage of the water containing phases \cite{richardson.1993,rossler.2006}.
%This causes them to melt into a gel-like structure, hence its early name.
Nowadays, this is less of an issue due to significant progress in the sensitivity of the detectors of electron microscopes.
However, to obtain enough signal for some detectors (e.g. \ac{EDX} or \ac{EBSD}), intensive electron beam exposure might still be required, which can cause visible damage \cite{rossler.2006, rossen.2017}.



\FloatBarrier
\section{Laser Ablation - Inductively Coupled Plasma - Mass Spectrometry (LA-ICP-MS)}

A \ac{LAICPMS} consists of a laser to ablate the specimen surface, a microscope to aim the laser at the measurement side, a \ac{ICP} to ionize the ablated flow of material within a carrier gas and a \ac{MS}.
This allows to detect minor and trace elements as well as isotopes in a spatially resolved manner.
Therefore, \ac{LAICPMS} is in wide use as an analytical method to determine the age of geological samples \cite{liu.2013} or for spatially resolved elemental mapping of heavy metals in biological tissue \cite{limbeck.2015}.

The laser ablation was done with a ns-Nd:YAG-laser at the quintupled wavelength of \SI{213}{\nano\meter} (\textit{ESI NWR 213}).
A quadrupole \ac{ICP}-\ac{MS} (\textit{NexION 300D} from \textit{Perkin Elmer}) was used for the elemental analysis.
The carrier gas chosen was helium due to its well-known advantages over other transport gases, such as its higher sensitivity for most \ac{ICP}-\ac{MS} instruments and smaller particle size distributions \cite{gunther.1999}.
The samples were placed in a two-volume sample cell and purged for at least \SI{10}{\min}.

\ac{LAICPMS} intrinsically results in a damaged surface, since it vaporizes material.
Figure~\ref{fig:laicpms-pre-ablation} illustrates this effect on a clinker surface.
It also indicates the the spot size (6 $\times$ \SI{6}{\micro\meter}) and the rough interaction depth, which is approximately \SI{7}{\micro\meter}.
On the positive side, this damage can be use for cross-method alignment purposes.

\begin{figure}[h!]
    \centering

	\ifdefined\PRINTABLE
		\includegraphics[height=.35\textheight]{la-icp-ms/pre-ablation.pdf}
	\else
		\animategraphics[palindrome,autoplay,height=.35\textheight]{4}{la-icp-ms/pre-ablation-anim/Animation}{00}{11}
	\fi
    \caption{\ac{SEM}-image of the pre-ablation and measurement area of a \ac{LAICPMS}-mapping. \cite{kleiner.2021d,kleiner.2021e}}
    \label{fig:laicpms-pre-ablation}
\end{figure}

It has to be noted, that specimen surface was pre-ablated using a coarser resolution and a larger area (see Figure~\ref{fig:laicpms-pre-ablation}) to remove surface contaminations.
However, this has to be taken into account to correctly align the final ablation result to the \ac{SEM}-image.

\FloatBarrier
\section{Pore size measurement}

In this chapter, methods used for pore size measurements are discussed.

\FloatBarrier
\subsection{Mercury Intrusion Porosimetry (MIP)}
\label{sec:method-mip}

In this thesis, an \textit{Autopore IV 9500} from \textit{Micromeritics GmbH} was used for the \ac{MIP} investigations.
It requires small pieces of dried specimen of about 5\,$\times$\,5\,$\times$\,\SI{5}{\cubic\micro\meter}.
It stepwise covers an pressure range of 0.003 to \SI{206}{\mega\pascal}, which roughly translates to 0.002 to \SI{200}{\micro\meter}.
However, values above {\color{red} ~\SI{20}{\micro\meter} should be ignored, since these are mostly cracks}.% geschätzt aus bettinas arbeit

\FloatBarrier
\subsection{Low Temperature - Differential Scanning Calorimetry (LT-DSC)}
\label{fig:ltdsc}

To obtain data about the capillary and gel porosity, the \textit{DSC Polyma 214} from \textit{Netzsch}, connected to the cooling system \textit{Intercooler IC70} was used.
The specimens were sliced to fit the small aluminum crucibles with a inner volume of $\approx$ \SI{60}{\micro\litre}.
These crucibles sealed air tight by cold-welding a lid under pressure and measured against an empty reference crucible.

% \begin{figure}[!htb]
%     \centering
%     \begin{tikzpicture}
% 		\pgfplotsset{every axis/.style={xmin=0, xmax=150,
% 		ymin=-65, ymax=25,
% 		width=0.8\linewidth,
% 		height=0.4\linewidth
% 		}}
%         \begin{axis}[
% 				legend style={at={(.97,.94)},anchor=north east},
% 				axis y line*=left,
% 				grid=major,
% 				xlabel={time $t$ in \si{\min}},
% 				ylabel={temperature $\vartheta$ in \si{\celsius}},
% 				xtick={0,20,...,150},
% 				ytick={-60,-40,...,20},
% 				grid style={line width=.1pt, draw=gray!10}
% 				%axis lines=middle
%             ]
% 			\addplot[mark=none, gray!40] coordinates {(0,0) (150,0)};
% 			\addplot [green, dashed,mark=.] coordinates {
% 				(0,20) (20,-20) (22,-20) (40,20) (48,-20) (49,-20) (50,-5) (60,0) (75,0)};
% 				\addplot [green, dashed,mark=.] coordinates {
% 					(95,-40) (105,-60) (110,-60) (120,-40)};
% 				\addplot [green, dashed,mark=.] coordinates {
% 					(140,0) (150,20) };

% 			\addplot[mark=none, gray] coordinates {(75,30) (75,-70)};

% 			%\node[black,right] at (axis cs:8010){\small{$W=0$}};
% 			\node[right] at (axis cs:78,10)  {$\rightarrow$ measurement};
% 			\node[right] at (axis cs:5,-40) {$\rightarrow$ inducing of ice nuclei};

%         \end{axis}
% 		% 3rd ordinate right
% 		\begin{axis}[blue,
% 			axis y line*=right,
% 			axis x line=none,
% 			ylabel={pore radius $r_\text{P}$ in \si{\nano\meter}},
% 			xtick={0,20,...,150},
% 			ytick      ={ -40, -30, -20, -10,  -1},
% 			yticklabels={2.19,2.27,3.80,7.04,65.2},
% 			ylabel style=black,
% 		  ]
% 		  \pgfplotsset{every outer y axis line/.style={xshift=1.25cm}, every tick/.style={xshift=1.25cm}, every y tick label/.style={xshift=1.25cm} }
% 		  \addplot [blue, dashed] coordinates {
% 			   (75,0) (95,-40)};

%         \end{axis}
% 		% 2nd ordinate right
% 		\begin{axis}[red,
% 			axis y line*=right,
% 			axis x line=none,
% 			xtick={0,20,...,150},
% 			ytick      ={ -40, -30, -20, -10,  -1},
% 			yticklabels={1.49,1.76,2.30,3.91,33.0},
% 		  ]
% 		  \addplot [red, dashed] coordinates {
% 			(120,-40) (140,0)};
%         \end{axis}
%     \end{tikzpicture}
%     \caption{Temperature regime of a \ac{LTDSC}-measurement, based on \textcite{sun.2010} and Müller \cite{mueller.2021}. theoretical pore radii after \textcite{brun.1977} - Red axis:  on melting (Equation \ref{eq:brun_melt}); blue axis: pore radius on freezing (Equation \ref{eq:brun_freeze}). The top value is the pore radius at \SI{-1}{\celsius}.}
%     \label{fig:ltdc_regime}
% \end{figure}

\begin{figure}[!htb]
    \centering
    \begin{tikzpicture}
		\pgfplotsset{every axis/.style={xmin=0, xmax=205,
		ymin=-65, ymax=25,
		width=0.8\linewidth,
		height=0.4\linewidth
		}}
        \begin{axis}[
				legend style={at={(.97,.94)},anchor=north east},
				axis y line*=left,
				grid=major,
				xlabel={time $t$ in \si{\min}},
				ylabel={temperature $\vartheta$ in \si{\celsius}},
				xtick={0,20,...,200},
				ytick={-60,-40,...,20},
				grid style={line width=.1pt, draw=gray!10}
				%axis lines=middle
            ]
			\addplot[mark=none, gray!40] coordinates {(0,0) (205,0)};
			\addplot [green, dashed,mark=.] coordinates {
				(0,20) (25,-30) (26,-30) (39,-5) (47,-1) (48,-1) (60,-25) (61,-25) (84,20) (109,-30) (110,-30) (123,-5) (131,-1) (132,-1) (162,-60) (163,-60) (203,20)};
				\addplot [blue,mark=.] coordinates {
					(48,-1) (60,-25)};
				\addplot [red,mark=.] coordinates {
					(61,-25) (73.8,0) };

				\addplot [blue,mark=.] coordinates {
					(132,-1) (152,-40)};
				\addplot [red,mark=.] coordinates {
					(173,-40) (193,0) };

			\addplot[mark=none, gray] coordinates {(48,30) (48,-70)};
			\addplot[mark=none, gray] coordinates {(73.8,30) (73.8,-70)};
			\addplot[mark=none, gray] coordinates {(132,30) (132,-70)};
			\addplot[mark=none, gray] coordinates {(193,30) (193,-70)};

			%\node[black,right] at (axis cs:8010){\small{$W=0$}};
			\node[right] at (axis cs:56,10)  {P1}; %$\rightarrow$
			\node[right] at (axis cs:155,10)  {P2}; %$\rightarrow$
			\node[right, blue] at (axis cs:18,-10) {\ding{100}1};
			\node[right, blue] at (axis cs:102,-10) {\ding{100}2};

        \end{axis}
		% 3rd ordinate right
		\begin{axis}[blue,
			axis y line*=right,
			axis x line=none,
			ylabel={pore radius $r_\text{P}$ in \si{\nano\meter}},
			xtick={0,20,...,150},
			ytick      ={ -40, -30, -20, -10,  -1},
			yticklabels={2.19,2.27,3.80,7.04,65.2},
			ylabel style=black,
		  ]
		  \pgfplotsset{every outer y axis line/.style={xshift=1.25cm}, every tick/.style={xshift=1.25cm}, every y tick label/.style={xshift=1.25cm} }
		  %\addplot [blue, dashed] coordinates {
			%   (75,0) (95,-40)};

        \end{axis}
		% 2nd ordinate right
		\begin{axis}[red,
			axis y line*=right,
			axis x line=none,
			xtick={0,20,...,150},
			ytick      ={ -40, -30, -20, -10,  -1},
			yticklabels={1.49,1.76,2.30,3.91,33.0},
		  ]
		  %\addplot [red, dashed] coordinates {
			%(120,-40) (140,0)};
        \end{axis}
    \end{tikzpicture}
    \caption{
		Temperature regime of a \ac{LTDSC}-measurement, based on \textcite{sun.2010} and \textcite{mueller.2021} with two measurement zones P1 and P2. {\color{blue}\ding{100}1}, {\color{blue}\ding{100}2} = zones are intended introduce ice nuclei. Theoretical pore radii are calculated after \textcite{brun.1977} - {\color{blue}Red axis}:  on melting (Equation \ref{eq:brun_melt}); blue axis: pore radius on freezing (Equation \ref{eq:brun_freeze}). The {\color{blue}65.2} and {\color{red}\SI{33.0}{\nano\meter}} are the pore radii at \SI{-1}{\celsius}.}
    \label{fig:ltdc_regime}
\end{figure}

The temperature regime used is shown in Figure~\ref{fig:ltdc_regime}.
The zones indicated with {\color{blue}\ding{100}1} and {\color{blue}\ding{100}2} are intended to create ice nuclei and therefore to avoid supercooling.
The measurement are done within the zones P1 (down to \SI{-25}{\celsius}) and P2 (down to \SI{-60}{\celsius}) at a cooling and heating rate of $\pm$\,\SI{2}{\kelvin\per\min}, which is the fastest usable heating rate according to \textcite{neffati.2001}.
These two measurement zones allow to distinguish the ice mass of pores with radii smaller and larger than $\approx$ \SI{3}{\nano\meter}.

The melting curve allows to process the cumulative ice mass.
Do to this, \textcite{chen.1988} and \textcite{ishikiriyama.1995} proposed the following.

% \begin{equation}
% 	\Delta T = T_\text{melt} - T
% 	\label{eq:ltdsc_delta_t}
% \end{equation}

% Randall to process the melting enthalpy
%\cite{randall} https://reader.library.cornell.edu/docviewer/digital?id=chla2944761_2176#page/19/mode/1up
% \Delta H^{0}_{m} =  \SI{334.1}{\joule\per\gram}
During supercooling the enthalpy of the fusion of ice ($H_\text{fus}(\Delta T)$ in \si{\joule\per\gram}) is changing slightly.
Therefore, at first this enthalpy has to be approximated for every $\Delta T$ using Equation~\ref{eq:ltdsc_h} by \textcite{chen.1988} and \textcite{ishikiriyama.1995} based on data from \textcite{randall.1930}:
\begin{equation}
	\Delta H_\text{fus}(\Delta T) = 334.1 + 2.119 \cdot \Delta T - 0.000783 \cdot \Delta T^2
	\label{eq:ltdsc_h}
\end{equation}

Then the heat flow data $\dot{Q}_{m}(\vartheta)$ caused by the phase transition has to be prepared by subtracting a baseline $Q_\text{B}(\vartheta)$ from the raw  data $Q_\text{raw}(\vartheta)$.
\begin{equation}
	{Q}_{m}(\vartheta) = Q_\text{raw}(\vartheta) -Q_\text{B}(\vartheta)
	\label{eq:mass_heat_flow}
\end{equation}

{\color{red}
The heat flow of the baseline is caused by the heat capacity of the whole specimen (binder + water/ice).
The heat capacity of the whole sample $mC_{p}$ is a combination of the heat capacities of the solid body ($C_{p\text{S}}$), the pore liquid ($C_{p\text{L}}$) and the ice crystals ($C_{p\text{I}}$) phase.
\begin{equation}
	mC_{p} = m_\text{S}C_{p} + m_\text{L}C_{p\text{L}} + m_\text{C}C_{p\text{I}}
	\label{eq:mass_heat_capacity}
\end{equation}

equation~30 and 31 from \textcite{sun.2010}:
\begin{equation}
	Q_\text{B}(\vartheta) = mC_{p} \cdot \frac{d\vartheta}{dt}
	\label{eq:ltdsc_baseline}
\end{equation}


In contrast, \textcite{wu.2014} and \textcite{johannesson.2010} used a dry reference sample to mitigate the heat capacity of the concrete binder.
%The baseline is calculated by a linear regression of the data points before and after the freezing peak as in {\color{red} \textsc{Unbehau} \cite{unbehau.2025}}.
However, this is a fairly subjective method, since the temperature (or time) ranges have to be chosen manually.


The calculation of the ice mass is influenced by the way, this baseline of the heat flow $Q_\text{B}$ is determined.
\begin{itemize}
	\item \textcite{sun.2010}, \textcite{mueller.2021} and \textcite{unbehau.2025} used a linear regression of the heat flow data utilizing manually selected temperature ranges before and after the freezing peak.
	\begin{itemize}
		\item \textcite{mueller.2021} analyzed only a small temperature change from \SI{-20}{\celsius} to \SI{+20}{\celsius} and therefore considered this method as sufficiently accurate.
		In both cases, the initial data were ignored, since they were .
	\end{itemize}
	\item 'J' baseline correction - \textcite{johannesson.2010} (large samples, slow heat rate (\SI{0.08}{\kelvin\per\minute}))
	\item 'C' baseline correction - \textcite{wu.2014} based on heat capacities of ice and water based on \cite{sun.2010}
\end{itemize}

ice density from \textcite{mueller.2021} calculated from data in \cite{eisenberg.2005}:

\begin{equation}
	\rho_\text{ice} = 0.9167 - \num{2.053e-4} \cdot \vartheta - \num{1.357e-6} \cdot \vartheta^2
	\label{eq:ice_density}
\end{equation}
}


% {\color{red}
% A automatic method was introduced to processes a baseline.
% \begin{itemize}
% 	\item The heat flow signal is denoised by calculating the mean value within a \SI{30}{\second} window of the dataset.
% 	\item For each value, the discrete difference to the previous value is calculated.
% 	\item The resulting signal is then denoised a second time as described in the first step.
% 	\item These steps increase the signal-to-noise ratio and allow to select the background and the signal.
% 	\item Raw values which fall between \num{-1} and \SI{4}{\nano\watt\per\milli\gram} after this processing, will then be used to calculate a linear regression which will then be used as baseline.
% \end{itemize}
% The detailed procedure ispublished in the function
% \texttt{auto\textunderscore select\textunderscore baseline} within the \texttt{ltdsc\textunderscore measurement}-class \cite{kleiner.2025}.
% } %\url{https://github.com/kleinerELM/Kalorimeterauswertung}
% % TODO: bild?

After the baseline correction, the total ice mass can be calculated by integrating the heat flow curve over the melting peak.
%\dot{Q}/m_\text{specimen} = DSC / Wärmestrom pro Probenmasse
\begin{equation}
	m_\text{ice} = \int_{\vartheta_\text{min}}^{\vartheta_\text{max}} \dot{Q}_{m}(\vartheta) d\vartheta  / \Delta H_\text{fus}
	\label{eq:ltdsc_ice}
\end{equation}

The ice mass of the gel pores can then be estimated by subtracting the ice mass of the capillary pores (P1-zone) from the total ice mass (P2-zone):

\begin{equation}
	m_\text{ice, gel} = m_\text{ice, P2} - m_\text{ice, P1}
	\label{eq:ltdsc_ice2}
\end{equation}
